% Migration IDs: P-02-001 through P-02-017 in SOURCE_MIGRATION_MANIFEST.csv.
\section{Literature and problem positioning}
\label{sec:literature}
\label{sec:literature-boundary}

The relevant literatures use similar terms--reduction, pruning, compression,
coverage, and libraries--for decisions that preserve different objects. We
therefore organize the comparison by research problem and preservation target.
Established work in operations research, stochastic control, and combinatorial
optimization comes first; recent agent-library and automated-strategy systems
are considered afterward as application-level comparators. This ordering is
important because the paper's covering algorithms are inherited from a
classical problem class, whereas its modeling contribution begins with the
dynamic semantics of a retained financial strategy library.

\subsection{Partial-information state reduction and semi-Markov control}

Belief-state reduction is foundational in partially observed control.
\citet{SmallwoodSondik1973} formulate finite-horizon control on the belief
simplex and establish the familiar piecewise-linear convex value
representation; \citet{Sondik1978} develops the discounted infinite-horizon
counterpart. Random decision durations do not remove the sufficient-statistic
logic: \citet{White1976} gives a finite-horizon, partially observed semi-Markov
dynamic program on a sufficient statistic. Financial switching and
regime-switching models likewise use filtered information to choose among a
fixed menu of operating modes or positions
\citep{Lakner1998,LiNystromOlofsson2015,DaiZhangZhu2010}.

Those results supply the control architecture used here, not a novelty claim.
The present reduction concerns a different component of state. Conditional on
the filtered belief and any active-project variables, it asks which summary of
an endogenous strategy library preserves both current operating choices and
the modeled law of future strategy generation. The answer under the declared
finite factorization is the operating frontier together with the generative
closure. Thus the frontier--closure pair summarizes the \emph{library}
component of the productive state; it neither replaces the belief nor proposes
a new general sufficient statistic for partially observed or semi-Markov
control.

Two standard reductions sharpen the distinction. Incremental pruning removes
alpha vectors that are redundant in a fixed partially observable dynamic-
programming backup \citep{CassandraLittmanZhang1997}. Its preserved object is
the pointwise value representation generated by that backup. A pruned vector
does not also act as a parent carrying modules for later policy generation.
State aggregation and model minimization instead identify states that are
behaviorally equivalent for a specified comparison class
\citep{GivanDeanGreig2003}. Their decision unit is a state or model component,
not a retained strategy entry. Innovation-safe compression therefore differs
from both alpha-vector pruning and state aggregation even though all three use
exact equality to justify a smaller representation.

\subsection{Policy, action, and control-library compression}

The closest established library-selection problem is static control-library
design. \citet{DeyEtAl2012} select and order compact maneuver or trajectory
libraries using submodular-sequence structure, trading library size and search
effort against expected task performance. The preserved or optimized object is
static task usefulness and search efficiency. In the present problem, a
retained entry can additionally alter the set of candidates that can be
generated later, so static control coverage does not certify generative
safety.

\citet{MuttiDelColRestelli2022} provide the closest policy-compression
comparison. They construct reward-free representatives of a policy space using
coverage of induced state--action distributions and formulate the selection as
set cover. Their preserved object is reward-independent distributional
representativeness for downstream reinforcement-learning tasks. Here the
feasible choices are sublibraries of a named source, and feasibility requires
exact equality of a regime- or belief-indexed operating frontier and a
parent-conditioned generative closure. Reward-free policy compression can
therefore retain a representative for broad downstream usefulness without
preserving the same objects as innovation-safe compression.

A directionally different problem expands rather than compresses the available
action set. In such models the performance object may be learning or regret in
a growing menu. Here generation is conditioned on modules carried by retained
parents, and the decision is which existing entries can be deleted without
changing either current opportunities or attainable descendants. The present
problem is therefore distinct from generative action-set expansion: it concerns
source-relative closure-preserving deletion rather than costly action creation.

\subsection{Weighted covering and combinatorial representation}

Set cover is the classical combinatorial core once a finite requirement
universe has been specified. Karp's set-covering decision problem is among the
foundational NP-complete problems \citep{Karp1972}. For positive column costs,
\citet{Chvatal1979} gives the standard binary weighted-cover formulation and
the minimum-cost-per-newly-covered-requirement greedy rule with its harmonic
approximation bound. These are established results. The paper neither
introduces weighted set cover nor claims a new generic greedy algorithm.

The contribution is the direction of the reduction. The dynamic semantic model
first determines what must be preserved: one tagged class records exact source
frontier attainment and another records source modules. Under identity closure,
a retained source sublibrary is innovation-safe if and only if it covers this
tagged requirement universe, after fixing the mandatory inactive strategy.
The tagged-cover equivalence is therefore a computational subclass
\emph{induced by} the dynamic frontier--closure model. It is not a semantic
interpretation imposed on an arbitrary set-cover instance, and it does not
extend automatically to a nonidentity closure with complementarities among
raw modules.

Computational benchmarking has a separate methodological literature.
Performance profiles compare solvers only after defining a common problem set
and success measure \citep{DolanMore2002}; feature-balanced benchmark
construction, as in MIPLIB 2017, is designed to avoid overrepresenting a narrow
instance family \citep{GleixnerEtAl2021}. These principles expose limits of the
present locked experiment rather than repairing it after the fact. Its methods
have unequal success targets, some infeasible generator cells, and a structured
block dominated by complete preprocessing. A future registered benchmark must
therefore balance realized residual structure and separate feasibility,
quality-qualified, and exact-completion targets before using comparative
profiles.

Adjacent covering models optimize different objects. Budgeted maximum coverage
and monotone submodular maximization seek the largest attainable value under a
cardinality or resource constraint \citep{NemhauserWolseyFisher1978}.
Submodular cover minimizes cost to reach the full value of a normalized,
monotone, integer-valued submodular function \citep{Wolsey1982}. Adaptive
submodularity concerns sequential choices whose outcomes are revealed over
time and requires adaptive monotonicity and diminishing returns
\citep{GolovinKrause2011}. Exact innovation-safe compression instead minimizes
positive active-strategy burden subject to complete coverage of fixed tagged
requirements. The maximum-coverage, submodular-cover, and adaptive guarantees
do not transfer unless their distinct assumptions are separately proved.

\subsection{Resource-constrained and metareasoning decisions}

Operations research has long treated information, search, computation, and
research budgets as scarce resources. \citet{Weitzman1979} optimizes costly
sequential search among alternatives. \citet{BakerEtAl1976} allocate an
organizational research-and-development budget across competing activities,
and \citet{RussellWefald1991} value computational actions by their effect on a
subsequent decision. These literatures motivate the economic importance of
research resources without supplying a preservation condition for a maintained
strategy library.

The distinction is between choosing what to investigate and choosing which
representation to retain. Exact innovation-safe compression is a one-time
outer decision that minimizes additive burden while holding the source
frontier and source closure fixed. The paper's capacity- and price-constrained
variants relax that equality and ask which productive state is worth buying;
they are closer to search, metareasoning, and research-budget allocation.
Neither variant turns held-out financial performance into a retention score.
This separation also explains why a locally safe deletion certificate answers
a feasibility question, whereas global minimum-burden retention answers a
resource-allocation question.

\subsection{Maintained generative libraries}

Library learning and program synthesis establish that reusable components can
shape later construction. For example, \citet{BowersEtAl2023} synthesize
libraries of reusable abstractions rather than treating candidate programs as
independent outputs. Its decision unit is internal program structure, and its
preserved object is reusable abstraction within synthesis. The present problem
instead selects a sublibrary of whole strategies that must jointly preserve a
belief-indexed operating envelope and parent-conditioned generation.
Innovation-safe compression is therefore not a substitute for procedural
compression, provenance, repair, versioning, or rollback controls.

\subsection{Automated financial strategy construction and governance}

Automated strategy construction predates contemporary agent systems.
\citet{AllenKarjalainen1999} use genetic search to construct technical trading
rules, and \citet{MoodySaffell2001} learn trading policies by direct
reinforcement. These studies optimize or evaluate generated strategies; they
do not pose minimum-burden retention of a persistent parent library. Likewise,
the model-risk literature studies the consequences of model uncertainty and
robust risk measurement \citep{Cont2006,GlassermanXu2014}. Supervisory guidance
documents expectations for model inventories, validation, monitoring, and
governance \citep{FederalReserveOCC2011}. This supports the existence of
organizational burdens, not the paper's optimization model or empirical
findings.

The finance literature also gives a direct warning about the interpretation of
large rule catalogs. \citet{SullivanTimmermannWhite1999} evaluate technical-rule
performance while accounting for the universe from which the rules were
selected, and \citet{White2000} develops a general reality-check framework for
data snooping. Multiple-testing concerns are especially acute in large return-
predictor collections \citep{HarveyLiuZhu2016}, while post-publication evidence
shows that reported anomaly returns can attenuate out of sample
\citep{McLeanPontiff2016}. The article's ex post maximum over enabled descendants
is not a corrected performance test under any of these frameworks. It is used
only to witness the mechanical implication of preserving or losing the declared
catalog, and no significance or investment claim is attached to it.

An accepted recent financial comparator, ReCAP, maintains regime-adaptive
portfolio policies \citep{PanEtAl2026ReCAP}. Its optimized object is
regime-specific portfolio performance, not minimum-burden preservation of a
source frontier and generative closure. The present paper begins after a finite
generation--verification--admission law has been declared and asks which source
entries can be removed without changing those two objects. Its financial
audits assess retention burden and semantic preservation, not improved returns,
forecasting ability, alpha, or deployable performance.

\medskip\noindent\textbf{Claim boundary.}
The paper claims novelty neither for state, value, policy, action-set,
control-library, or internal procedural compression nor for weighted cover,
metareasoning, maintained skill libraries, or automated strategy generation.
The narrower contribution is that, under the declared finite dynamic model,
the operating frontier and generative closure jointly define the object that
safe retention must preserve; under identity closure, those semantics induce
an exact tagged weighted-cover subclass and source-relative certificates.
None of the verified comparisons reviewed here establishes that joint
formulation. This is a bounded audit conclusion, not a priority claim over all
related library work.
